\chapter{METODE DAN RANCANGAN PENELITIAN}

\section{Deskripsi Umum Penelitian}

Penelitian ini menggunakan metode eksperimen terkontrol untuk mengevaluasi pengaruh kuantisasi terhadap kualitas dan efisiensi beberapa Small Language Model pada klasifikasi sentimen Bahasa Indonesia. Setiap varian kuantisasi dibandingkan dengan baseline presisi tinggi dari model yang sama. Dengan demikian, analisis utama tidak hanya membandingkan skor absolut antarmodel, tetapi juga perubahan relatif dan transisi prediksi pada sampel yang sama.

\section{Metodologi}

Tahapan penelitian dilakukan secara sistematis sebagai berikut:

Tahapan penelitian dilakukan secara sistematis sebagai berikut:

\begin{enumerate}
    \item \textbf{Tahap 1:} Melakukan penelusuran penelitian terkait dan menetapkan kriteria pemilihan model serta dataset.
    \item \textbf{Tahap 2:} Melakukan pilot kompatibilitas model, tokenizer, chat template, parser, dan kebutuhan sumber daya.
    \item \textbf{Tahap 3:} Mengonversi checkpoint model ke GGUF presisi tinggi dan menghasilkan varian kuantisasi standar secara konsisten.
    \item \textbf{Tahap 4:} Menjalankan inferensi pada dataset klasifikasi sentimen dengan prompt, decoding, dan runtime yang dikendalikan.
    \item \textbf{Tahap 5:} Mengukur kualitas, kepatuhan format, transisi prediksi, ukuran berkas, memori, latensi, dan throughput.
    \item \textbf{Tahap 6:} Menganalisis degradasi relatif dan menentukan konfigurasi pada Pareto frontier kualitas--efisiensi.
\end{enumerate}

\section{Rancangan Eksperimen}

\subsection{Alur Eksperimen}

Rancangan eksperimen terdiri dari beberapa komponen utama:

\begin{itemize}
    \item \textbf{Checkpoint model:} empat model utama yang dipilih berdasarkan hasil pilot.
    \item \textbf{Konversi dan kuantisasi:} pembentukan baseline F16/BF16 serta varian Q8\_0, Q6\_K, Q5\_K\_M, Q4\_K\_M, dan IQ4\_XS atau IQ4\_NL.
    \item \textbf{Runtime inferensi:} satu versi atau commit \texttt{llama.cpp} dengan konfigurasi yang sama.
    \item \textbf{Evaluator:} modul prompt, parser label, penghitung metrik, penyimpan keluaran per sampel, dan pengukur performa sistem.
\end{itemize}

% TODO: Ganti dengan diagram arsitektur sistem Anda.
% \begin{figure}[H]
%   \centering
%   \includegraphics[width=0.9\textwidth]{figures/arsitektur-sistem}
%   \caption{Diagram Arsitektur Sistem}
%   \label{fig:arsitektur}
% \end{figure}

\subsection{Bagan Alir Metodologi}

% TODO: Tambahkan bagan alir (flowchart) metodologi penelitian.

% TODO: Ganti dengan bagan alir metodologi Anda.
% \begin{figure}[H]
%   \centering
%   \includegraphics[width=0.8\textwidth]{figures/bagan-alir}
%   \caption{Bagan Alir Metodologi Penelitian}
%   \label{fig:bagan-alir}
% \end{figure}

\section{Lingkungan Penelitian}

Lingkungan penelitian yang digunakan dalam eksperimen utama adalah:

\begin{itemize}
    \item \textbf{Perangkat lunak:} satu versi atau commit \texttt{llama.cpp}, tokenizer dan evaluator yang terdokumentasi.
    \item \textbf{Perangkat keras:} AMD Ryzen 3 5300U, RAM 9 GB, tanpa GPU.
    \item \textbf{Konfigurasi:} jumlah thread, panjang konteks, panjang keluaran, seed, prompt, dan parameter decoding dikendalikan.
\end{itemize}

\section{Metode Pengukuran dan Analisis}

Pengukuran dilakukan dengan beberapa pendekatan:

\begin{enumerate}
    \item \textbf{Kualitas prediksi:} akurasi dan macro-F1 pada dataset klasifikasi sentimen.
    \item \textbf{Kepatuhan dan transisi:} valid-label rate, prediction agreement, serta kategori benar-ke-salah dan salah-ke-benar.
    \item \textbf{Efisiensi:} ukuran berkas, peak RAM, waktu pemuatan, latensi, dan throughput.
    \item \textbf{Analisis:} selisih absolut, degradasi relatif terhadap baseline model yang sama, interval kepercayaan bootstrap berpasangan, dan Pareto frontier.
    \item \textbf{Reproduksibilitas:} pencatatan versi model dan runtime, hash berkas, konfigurasi, prompt, seed, dan hasil per sampel.
\end{enumerate}
